\chapter{Leukemia}
\index{Leukemia}

\section*{Definition and Overview}
Leukemia (referred to in British and local medical terminology as leukaemia) represents a heterogeneous group of malignant haematological neoplasms originating in the bone marrow. It is characterised by the clonal proliferation and accumulation of abnormal, immature white blood cells (blasts) that fail to differentiate normally. These malignant cells crowd out the bone marrow, progressively suppressing normal haematopoiesis and leading to severe cytopenias. Leukemia is broadly classified based on the speed of progression into acute and chronic forms, and further subdivided by the cell lineage involved into myeloid or lymphoid. The four primary subtypes are acute myeloid leukaemia (AML), acute lymphoblastic leukaemia (ALL), chronic myeloid leukaemia (CML), and chronic lymphocytic leukaemia (CLL).

\section*{Epidemiology}
Leukaemias account for approximately 3\% of all new cancer diagnoses globally. In many countries, leukaemia is a significant cause of oncology-related morbidity and mortality, with thousands of cases diagnosed annually. The epidemiology varies dramatically by subtype: ALL is predominantly a disease of childhood, representing the most common paediatric cancer, with a peak incidence between 2 and 5 years of age. In contrast, AML, CLL, and CML are primarily diseases of older adults, with the median age of diagnosis for CLL and AML being over 70 years. There is a slight male predominance across most subtypes. Environmental and genetic factors contribute to regional variation, though the incidence remains highest in developed nations, likely reflecting advanced diagnostic capabilities and ageing populations.

\section*{Aetiology and Pathophysiology}
Leukaemia arises from somatic genetic mutations in haematopoietic stem or progenitor cells. The primary pathogenetic mechanisms and risk factors include:
\begin{itemize}
    \item \textbf{Chromosomal translocations:} Reciprocal translocations that create novel fusion oncogenes, such as the Philadelphia chromosome t(9;22) creating the \textit{BCR-ABL1} fusion gene in CML and a subset of ALL, which results in constitutive tyrosine kinase activity.
    \item \textbf{Genetic predisposition:} Conditions such as Down syndrome (trisomy 21), Fanconi anaemia, and ataxia-telangiectasia are associated with a significantly increased risk of developing acute leukaemia.
    \item \textbf{Environmental exposures:} Exposure to high-dose ionising radiation and chemical carcinogens, particularly benzene and other organic solvents.
    \item \textbf{Prior oncological therapy:} Development of therapy-related AML (t-AML) following chemotherapy with alkylating agents or topoisomerase II inhibitors, or after therapeutic radiation.
    \item \textbf{Marrow replacement and cytopenia:} The clonal expansion of non-functional leukaemic blasts physically disrupts the bone marrow microenvironment, leading to erythropenia (causing anaemia), thrombocytopenia (causing bleeding), and functional neutropenia (causing infection susceptibility).
\end{itemize}

\section*{Clinical Features}
The clinical presentation of leukaemia reflects the failure of normal bone marrow function and the infiltration of malignant cells into various tissues. Common clinical features include:
\begin{itemize}
    \item \textbf{Anaemic syndrome:} Profound fatigue, generalized weakness, pallor, and exertional dyspnoea due to reduced red blood cell production.
    \item \textbf{Bleeding diathesis:} Easy bruising, petechiae (particularly on the lower extremities), epistaxis, menorrhagia, and bleeding gums resulting from thrombocytopenia.
    \item \textbf{Infection and fever:} Frequent, severe, or atypical infections (e.g. oral candidiasis, atypical pneumonia) and persistent fevers due to a lack of functional neutrophils.
    \item \textbf{Constitutional symptoms:} Unintentional weight loss, drenching night sweats, and anorexia, which are particularly prominent in advanced or acute leukaemias.
    \item \textbf{Organ infiltration:} Bone and joint pain due to marrow expansion; hepatomegaly and splenomegaly (frequently causing left upper quadrant pain or early satiety); and painless lymphadenopathy (common in ALL and CLL).
    \item \textbf{Leukostasis:} Symptoms of microvascular occlusion (headache, visual changes, dyspnoea, or confusion) occurring in patients with acute leukaemias who present with extremely high white blood cell counts (typically $>100 \times 10^9$/L).
\end{itemize}

\section*{Differential Diagnosis}
The symptoms of leukaemia overlap with many benign and malignant disorders. The primary differential diagnoses include:
\begin{itemize}
    \item \textbf{Leukaemoid reaction:} Severe neutrophilic leukocytosis (typically $>50 \times 10^9$/L) secondary to severe sepsis, physical stress, or solid tumours, distinguished by a high leukocyte alkaline phosphatase (LAP) score and absence of clonal blasts.
    \item \textbf{Myelodysplastic syndromes (MDS):} Clonal bone marrow disorders presenting with cytopenias, but characterized by dysplastic cells and a bone marrow blast count of less than 20\%.
    \item \textbf{Aplastic anaemia:} Hypocellular marrow leading to pancytopenia, distinguished from leukaemia by the complete absence of circulating blasts and lack of hepatosplenomegaly or lymphadenopathy.
    \item \textbf{Infectious mononucleosis:} Acute Epstein-Barr virus (EBV) or cytomegalovirus (CMV) infection presenting with fever, pharyngitis, lymphadenopathy, and atypical lymphocytosis on blood film, which can mimic lymphoid leukaemia.
    \item \textbf{Other bone marrow malignancies:} Multiple myeloma or metastatic solid tumours (such as breast or lung carcinoma) infiltrating the bone marrow space.
\end{itemize}

\section*{When to Seek Medical Care}
Patients presenting with unexplained, persistent fatigue, easy bruising, recurrent infections, or unexplained fevers should undergo prompt medical evaluation. A full blood count should be performed rapidly.

\textbf{Contact your local emergency services for:}
\begin{itemize}
    \item A fever of 38 degrees Celsius or higher in a patient undergoing chemotherapy (neutropenic sepsis).
    \item Severe, uncontrollable bleeding, or signs of internal haemorrhage (e.g. vomiting blood, black tarry stools).
    \item Sudden shortness of breath, chest pain, or severe headache.
    \item Signs of leukostasis, including sudden confusion, visual disturbances, or focal neurological deficits.
\end{itemize}

\section*{Diagnosis and Investigations}
A definitive diagnosis of leukaemia requires evaluation of peripheral blood and bone marrow:
\begin{itemize}
    \item \textbf{Full blood count (FBC) and blood film:} Demonstrates cytopenias and often shows marked leukocytosis with circulating blast cells or atypical lymphocytes.
    \item \textbf{Bone marrow aspirate and trephine biopsy:} Essential for morphological examination; a blast count of 20\% or greater in the marrow or blood is diagnostic of acute leukaemia.
    \item \textbf{Flow cytometry (Immunophenotyping):} Performed on blood or marrow to identify lineage-specific cluster of differentiation (CD) markers (e.g. CD19, CD20 for B-ALL; CD3, CD8 for T-ALL; CD13, CD33, myeloperoxidase for AML).
    \item \textbf{Cytogenetics (Karyotyping and FISH):} Detects chromosomal abnormalities which are vital for diagnosis and risk-stratification (e.g. t(15;17) in acute promyelocytic leukaemia; t(9;22) in CML).
    \item \textbf{Molecular genetics:} Polymerase chain reaction (PCR) or next-generation sequencing (NGS) to identify prognostic mutations such as \textit{FLT3}, \textit{NPM1}, or \textit{JAK2}.
    \item \textbf{Coagulation studies:} Essential to rule out disseminated intravascular coagulation (DIC), particularly if acute promyelocytic leukaemia is suspected.
\end{itemize}

\section*{Management}
Treatment is highly individualised and depends on the specific subtype, cytogenetic risk profile, patient age, and comorbidities:
\begin{itemize}
    \item \textbf{Induction chemotherapy:} Intensive multi-agent chemotherapy regimens aimed at rapidly achieving complete haematological remission.
    \item \textbf{Consolidation and maintenance therapy:} Post-remission therapy (chemotherapy or immunotherapy) to eradicate minimal residual disease (MRD) and prevent relapse.
    \item \textbf{Targeted therapies:} Small-molecule inhibitors such as tyrosine kinase inhibitors (e.g. imatinib, dasatinib) for CML and Philadelphia-positive ALL, or FLT3 inhibitors (e.g. midostaurin) for AML.
    \item \textbf{Immunotherapy:} Monoclonal antibodies (e.g. rituximab for CLL), bispecific T-cell engagers (e.g. blinatumomab), or chimeric antigen receptor (CAR) T-cell therapy for relapsed lymphoblastic malignancies.
    \item \textbf{Allogeneic haematopoietic stem cell transplant (HSCT):} Reserved for high-risk acute leukaemias or relapsed disease to provide a graft-versus-leukaemia effect.
    \item \textbf{Supportive care:} Prophylactic antimicrobials, transfusion of irradiated blood products, and management of metabolic complications like tumour lysis syndrome.
\end{itemize}

\section*{Complications}
The complications of leukaemia and its treatments can be life-threatening and include:
\begin{itemize}
    \item \textbf{Neutropenic sepsis:} Severe systemic infection in the absence of adequate white blood cells, requiring emergency broad-spectrum intravenous antibiotics.
    \item \textbf{Tumour lysis syndrome (TLS):} Metabolic derangement (hyperkalaemia, hyperuricaemia, hyperphosphataemia, hypocalcaemia) caused by rapid cell breakdown, which can lead to acute renal failure.
    \item \textbf{Haemorrhage:} Severe bleeding, particularly intracranial or pulmonary, secondary to profound thrombocytopenia or DIC.
    \item \textbf{Graft-versus-host disease (GVHD):} An immunological reaction where donor immune cells attack host tissues following an allogeneic HSCT.
    \item \textbf{Long-term treatment sequelae:} Secondary malignancies, cardiotoxicity, infertility, and cognitive changes (chemo-brain).
\end{itemize}

\section*{Prognosis and Outcomes}
The prognosis of leukaemia has improved dramatically over the past few decades but remains highly variable. Paediatric ALL has an outstanding prognosis, with 5-year survival rates exceeding 90\%. Chronic lymphocytic leukaemia is often slow-growing, and many elderly patients have a normal life expectancy without ever requiring active treatment. In contrast, AML in older adults carries a guarded prognosis, with 5-year survival rates often below 20\% due to treatment-refractory disease and high rates of relapse. Chronic myeloid leukaemia has a near-normal life expectancy due to the success of lifelong tyrosine kinase inhibitor therapy.

\section*{Prevention and Self-Care}
While most cases of leukaemia cannot be prevented, patients can take specific self-care steps to manage the disease and treatments:
\begin{itemize}
    \item \textbf{Infection avoidance:} Avoid crowds, sick contacts, and unpasteurised or raw foods during periods of neutropenia (neutropenic diet).
    \item \textbf{Meticulous hygiene:} Wash hands frequently, use a soft toothbrush to prevent gum trauma, and maintain excellent skin integrity.
    \item \textbf{Vaccination adherence:} Ensure up-to-date immunisation with inactivated vaccines (such as the annual influenza vaccine) while strictly avoiding live vaccines.
    \item \textbf{Hydration:} Drink plenty of fluids to assist in flushing chemotherapy metabolites and preventing renal injury.
    \item \textbf{Psychosocial support:} Engage with oncology social workers, psychologists, and support organisations to manage the significant emotional burden of a leukaemia diagnosis.
\end{itemize}

\section*{Sources and Further Reading}
The following reputable organisations provide further information on leukaemia:
\begin{itemize}
    \item Leukaemia Foundation (Australia). \textit{Support services and patient guides for blood cancers.} https://www.leukaemia.org.au/
    \item Cancer Council Australia. \textit{Information on leukaemia diagnosis, staging, and support.} https://www.cancer.org.au/
    \item Leukemia \& Lymphoma Society (LLS). \textit{Comprehensive educational resources on blood cancers.} https://www.lls.org/
    \item NHS. \textit{Overview and treatment of acute and chronic leukaemia.} https://www.nhs.uk/
    \item Mayo Clinic. \textit{Clinical information on leukaemia types, symptoms, and trials.} https://www.mayoclinic.org/
\end{itemize}
